\section{什么是 LLM Agent}

\subsection{Agent 的系统组成}

\begin{figure}[H]
\centering
\includegraphics[page=4,width=0.9\textwidth]{slides.pdf}
\caption{课程用 Memory、Tool Use、Retrieval、Reasoning/Planning 和 Environment Feedback 概括 Agent 的标准组件。}
\end{figure}

课程给出的图景很清晰：Agent 的中心仍是 LLM，但 LLM 不再孤立运行，而是通过多个接口与外部系统耦合。

\begin{itemize}
    \item \textbf{Memory}：保存任务上下文、用户偏好、历史轨迹，避免每轮对话都从零开始。
    \item \textbf{Tool use}：让模型调用搜索、代码执行、API、数据库或浏览器操作。
    \item \textbf{Retrieval}：从文档和知识库中拉取外部证据，避免全靠参数记忆硬答。
    \item \textbf{Reasoning \& Planning}：把复杂目标拆成若干步骤，而不是直接给出单轮输出。
    \item \textbf{Feedback / Environment}：让 Agent 根据执行结果回看自己，形成 trial-and-error 的闭环。
\end{itemize}

\begin{knowledgebox}{为什么要保留这些模块边界}
如果把记忆、检索、工具调用和规划全都塞回一个超长 prompt，系统将很难观测、调试和优化。Agent 框架的重要价值就在于把这些能力显式拆开，使每个组件都可以被替换、约束和评估。
\end{knowledgebox}

\subsection{Agent 活跃的环境类型}

\begin{figure}[H]
\centering
\includegraphics[page=5,width=0.9\textwidth]{slides.pdf}
\caption{LLM Agent 不只存在于聊天界面，还可以进入网页、软件开发、企业工作流和机器人等多种环境。}
\end{figure}

课程特别强调 Agent 的环境多样性。它既可以工作在数字环境中，例如 IDE、浏览器、工作流平台，也可以延伸到具身环境，如机器人和多模态感知任务。环境越开放，Agent 越需要：
\begin{itemize}
    \item 在不完整信息下做局部决策；
    \item 从失败中恢复，而不是一次错误直接终止；
    \item 同时协调符号操作、自然语言交互和外部动作。
\end{itemize}

\subsection{本章小结}

LLM Agent 的定义不是``会调用工具的 LLM''这么简单，而是一个由模型、外部能力和反馈回路构成的任务执行系统。

